% TOP_PROBLEM: {"problemId": "problem.mumford-shah-conjecture-on-the-jump-set-of-planar-minimizers", "canonicalTitle": "Planar homogeneous Mumford–Shah local-structure conjecture", "releaseRank": 279, "primaryDomain": "analysis_pde", "primaryDomainLabel": "Analysis & PDE", "displayStatus": "open_with_solved_subcases", "releaseStatus": "open_with_solved_subcases"}
% !TeX program = lualatex
% ATTEMPT_STATUS: unresolved
\documentclass[11pt]{article}
\usepackage[margin=25mm]{geometry}
\usepackage{fontspec}
\setmainfont{Latin Modern Roman}
\usepackage{amsmath,amssymb,mathtools}
\usepackage{unicode-math}
\setmathfont{Latin Modern Math}
\usepackage{xurl,hyperref}
\hypersetup{colorlinks=true,urlcolor=blue}
\setcounter{secnumdepth}{0}
\setlength{\parindent}{0pt}
\setlength{\parskip}{5pt}
\setlength{\emergencystretch}{3em}
\begin{document}
\section*{\texorpdfstring{Planar homogeneous Mumford–Shah local-structure conjecture}{Planar homogeneous Mumford–Shah local-structure conjecture}}\label{279}
\subsection{Definitions and mathematical statement}
For $u\in W^{1,2}_{\mathrm{loc}}(\Omega\setminus K)$ and relatively closed countably rectifiable $K\subset\Omega\subset{\mathbb{R}}^2$, the homogeneous Mumford--Shah energy is
$E_0(u,K;U)=\int_{U\setminus K}|\nabla u|^2+\mathcal H^1(K\cap U)$.
Use reducedness and absolute compactly supported local minimality exactly as stated below, allowing all admissible changes of the crack set, without a separation constraint. The conjecture says that near each $x\in K$, a $C^1$ diffeomorphism fixing $x$ with derivative $I$ straightens $K$ to a diameter, a single radius ending at $x$, or three radii meeting at $120$ degrees. Thus only a regular crack, crack tip, or triple junction is allowed. The source's hypotheses on local finite energy and negligible removable crack pieces are part of the statement; this is not a classification of arbitrary stationary pairs.
\par\smallskip\textit{Formulation and concept references:} \href{https://arxiv.org/pdf/2308.14660v2}{[S1]}.

\subsection{Short English statement}
Must every reduced planar Mumford–Shah local minimizer have only smooth crack arcs, crack tips, and three-way junctions with 120-degree angles?
\subsection{Research attempt}
\paragraph{Scope, source check, and attack.}
This attempt was prepared by GPT-6 Astra Ultra on 27 September 2026.
The source [S1] distinguishes absolute minimizers from competitors
constrained to preserve separation, and from stationary pairs.
Its classification and regularity results use those distinctions.
We retain absolute compactly supported local minimality and reducedness.
The work derives necessary equations on already regular pieces,
checks the crack-tip coefficient and its borderline integrability,
and gives an explicit stationary junction which fails minimality.
It does not deduce regular pieces at every point from those equations.

Write $E(u,K;U)=\int_{U\setminus K}|\nabla u|^2+
\mathcal H^1(K\cap U)$. Changes of $K$ are allowed inside a compact
subset of $\Omega$, and no artificial separation condition is added.
The desired straightening has derivative identity after the comparison
cone is chosen with the appropriate orientation; a rotation is not
silently built into that derivative requirement.

\paragraph{The correct blow-up normalization.}
At $x_0\in K$ and scale $r>0$, set
\[
 K_r=(K-x_0)/r,\qquad
 u_r(y)=\frac{u(x_0+ry)-c_r}{\sqrt r}.
\]
For any constant $c_r$, the change of variables gives
\[
           E(u_r,K_r;B_R)=r^{-1}E(u,K;B_{rR}(x_0)).       \tag{1}
\]
Indeed $\nabla u_r=\sqrt r\,\nabla u(x_0+ry)$, while planar area
scales by $r^{-2}$ and crack length by $r^{-1}$.
The amplitude factor $r^{-1/2}$ is essential. Scaling $u$ only by
translation would not keep the two energy terms at the same order.

Constants on different complementary components can separate without
bound under this normalization. A bound on $\int|\nabla u_r|^2$ does
not therefore imply a global bound on $\int|u_r|^2$ after subtracting
just one constant. Any compactness argument must handle componentwise
normalization and the possible changing crack topology. Formula (1)
alone does not establish a compactness or classification theorem.

\paragraph{Necessary equations where smoothness is already known.}
Suppose a portion of $K$ is a smooth embedded arc and $u$ is smooth on
each side. Varying $u$ with $K$ fixed and support away from the arc gives
$\Delta u=0$ in each component. Allowing independent trace variations
on the two faces gives zero normal derivative on each face:
\[
                              \partial_\nu u^\pm=0.      \tag{2}
\]
These statements follow by integration by parts in the Dirichlet energy.
They are conditions on a regular portion, not a proof that every
portion is regular.

Choose a unit normal $n$ pointing from the minus side to the plus side,
and write the curvature vector of the arc as $\kappa n$.
Moving the arc by $h n$ enlarges the minus domain and shrinks the plus
domain. The boundary term from the bulk energy, using (2), is
$(|\nabla u^-|^2-|\nabla u^+|^2)h$; the length variation is
$-\kappa h$. Stationarity hence requires
\[
                         \kappa=|\nabla u^-|^2-|\nabla u^+|^2.
                                                               \tag{3}
\]
Changing the chosen normal reverses both sides. This convention check
prevents a curvature sign from being used inconsistently.

At a junction with three regular arcs and bounded bulk gradients,
the first variation under displacement of the junction has no
point-mass contribution from the bulk energy. The length contribution
forces the sum of the three outgoing unit tangents to be zero.
Three unit vectors sum to zero precisely when their pairwise angles
are $120$ degrees: for example
$|\tau_1+\tau_2|^2=|\tau_3|^2=1$ gives
$\tau_1\cdot\tau_2=-1/2$, and similarly for the other pairs.
The bounded-gradient regular-junction assumption is explicit here;
it cannot be used to rule out all possible singular junctions.

\paragraph{An explicit four-way stationary failure.}
Take the four terminals $(\pm a,\pm a)$ and connect each to the origin.
The resulting cross has total length $4\sqrt2\,a$.
Assign a different constant value of $u$ to each of its four sectors.
The bulk energy is zero; each edge is straight, and the four unit
tangents at the origin sum to zero. Thus the smooth first-variation
conditions alone do not reject this configuration.

Replace the cross inside the square by a Steiner tree with junctions
\[
                  z_\pm=(\pm(a-a/\sqrt3),0).
\]
Join the two left corners to $z_-$, the two right corners to $z_+$,
and join $z_-$ to $z_+$. Each of the four outer edges has length
$2a/\sqrt3$, and the middle edge has length $2a-2a/\sqrt3$.
The new total is
\[
                        2a(1+\sqrt3)<4\sqrt2\,a.        \tag{4}
\]
The inequality follows, for instance, by squaring
$1+\sqrt3<2\sqrt2$. The new tree has four complementary regions,
one adjoining each boundary arc between consecutive terminals.
Assign to each region the constant prescribed on that boundary arc.
This matches the original outside values and keeps the bulk energy zero.
Extend both configurations by the same exterior rays in a slightly
larger disk, so the modification is compactly supported. Equation (4)
strictly reduces the total energy.

This is an admissible topology-changing competitor. It shows directly
why stationarity cannot replace absolute minimality. It is not a
counterexample to the conjecture, since the original cross is not a
minimizer. Nor does it classify arbitrary junctions whose bulk energy
may be singular rather than zero.

\paragraph{The crack-tip coefficient from an energy-release calculation.}
Let $K=\{(x,0):x\le0\}$ and use polar angles $-\pi<\theta<\pi$.
Consider
\[
                         u(r,\theta)=A\sqrt r\sin(\theta/2).
                                                               \tag{5}
\]
It is harmonic off the slit and satisfies (2) on the two faces.
Direct differentiation gives
\[
 |\nabla u|^2=\frac{A^2}{4r},\qquad
 \int_{B_R\setminus K}|\nabla u|^2=\frac{\pi A^2}{2}R,
 \qquad E(u,K;B_R)=\left(1+\frac{\pi A^2}{2}\right)R.    \tag{6}
\]
Harmonicity and the face condition hold for every $A$, so they do not
fix the normalization required for tip stationarity.

For the energy convention without a factor $1/2$, the translation
stress flux in the positive horizontal direction is
\[
 J=\int_{\partial B_R}
      \bigl(|\nabla u|^2\nu_1-2\partial_\nu u\,\partial_1u\bigr)\,ds.
\]
Here $\partial_\nu u=A\sin(\theta/2)/(2\sqrt R)$ and
$\partial_1u=-A\sin(\theta/2)/(2\sqrt R)$.
The integrand reduces to $A^2/(4R)$, because
$\cos\theta+2\sin^2(\theta/2)=1$. Thus
\[
                                  J=\pi A^2/2.          \tag{7}
\]
The slit faces give no horizontal flux. Translating the endpoint to
extend the crack adds unit length per unit displacement, while the
bulk energy release is $J$. Stationarity under a localized tip
translation requires $1-J=0$, hence
\[
                                  |A|=\sqrt{2/\pi}.     \tag{8}
\]
This is a necessary normalization calculation. A proof that the
normalized crack-tip pair minimizes against all admissible competitors
requires more than this first variation; that fact belongs to the
deeper theory discussed in [S1] and is not reproved by (7).

\paragraph{The sharp integrability obstruction is visible in the model.}
For (5), $|\nabla u|=|A|/(2\sqrt r)$. For $q>0$,
\[
 \int_{B_R}|\nabla u|^q
    =2\pi(|A|/2)^q\int_0^R r^{1-q/2}\,dr.
\]
For $A\ne0$ this is finite exactly when $q<4$.
At $q=4$ it diverges logarithmically, but the distribution function is
\[
 |\{x\in B_R:|\nabla u(x)|>\lambda\}|
     =\pi\min\left\{R^2,\frac{A^4}{16\lambda^4}\right\}. \tag{9}
\]
Thus the gradient belongs to weak $L^4$, even though it does not belong
to strong $L^4$. This explains why a proposed route demanding strong
$L^4$ estimates for every minimizer would exclude a legitimate crack-tip
model. The source's sharper regularity discussion uses the weak-space
threshold; a generic higher-integrability estimate with some exponent
$2+\epsilon$ does not reach it.

\paragraph{Where the local classification remains missing.}
Equations (2)--(3), the junction force balance, and the model calculation
all start with geometric or analytic structure already available.
They do not show that every blow-up has one of the listed forms.
Even a subsequential model blow-up would require a quantitative
regularity theorem to produce a unique tangent configuration and the
stated $C^1$ straightening at the original point. Treating subsequential
Hausdorff convergence as a differentiable parametrization skips that
step. Reducedness also matters: adding a removable isolated crack
point changes the set without changing either energy term.

\paragraph{Verification and final status.}
The local checks verify the Steiner-tree lengths and angles, the polar
derivatives in (6)--(7), and the exponents in (1) and (9). No finite
simulation certifies absolute minimality against all crack changes.

\textbf{UNRESOLVED.} The checked results are necessary regular-piece
equations, a strict competitor excluding a four-way zero-energy cross,
and the crack-tip energy and weak-$L^4$ calculations. The first missing
step is exclusion of nonmodel blow-up behavior for arbitrary reduced
absolute minimizers. Specific next targets are compactness with
componentwise constants, a classification theorem with the correct
competitor class, and quantitative improvement from model closeness to
$C^1$ geometry. None is replaced here by stationarity alone.

\subsection{Sources}
\begin{itemize}
\item[S1]The regularity theory for the Mumford-Shah functional on the plane; arXiv2308.14660v2 January5,2025.
\url{https://arxiv.org/pdf/2308.14660v2}.
\textit{Formulation source.}
\end{itemize}
\end{document}
